\hwhat{Can a swarm's unforced fluctuations predict how it lines up behind an assigned goal? Each agent's 30-min statement embeddings were projected on the goal direction $\hat g$ (goal plus kickoff text). In three free $\to$ assigned pairs (\#11$\to$\#12, \#16$\to$\#17, \#37$\to$\#38, non-holdout), one field $\lambda$ was fitted to the mean shift and everything else predicted. Alignment rose in every pair, by 1.4--3.1 free-week SDs, outside linear response, so the variance test was untestable. \textbf{The core test failed:} free-week fluctuations did not predict who moved ($r=-0.37,-0.69,+0.20$; Stouffer $p=0.89$; synthetic power 95\%). Variance along $\hat g$ instead grew $\times$1.9--7.7 while transverse variance held, the opposite of a saturating tilt. Across 18 regime-I goal changes the jump fell with prior fluctuation (Spearman $-0.56$, $p=0.015$). Agents converge on a common alignment and switch on- and off-goal together. Holdout (\#22$\to$\#23, \#31$\to$\#32) frozen, not run.}

\hmodel{Model 11, mean-field O($n$) in $n=32$ whitened embedding dimensions. $\mathbf s_i$: agent $i$'s unit statement vectors; $y=\hat{\mathbf z}\cdot\hat g$, alignment with the goal field direction; $x_i$, its 30-min window mean, with free-week mean $\mu_i$, noise-corrected variance $\kappa_{2,i}$, third cumulant $\kappa_{3,i}$. The goal is a linear field $\lambda=\beta h$ along $\hat g$; $K$ is the cumulant generating function; $J_0$ the mean-field coupling; $g$ the loop gain; $C_F$ the free-week window covariance.}

\hmath{\vspace{-6pt}\begin{gather*}
P_A(\mathbf s)\propto P_F(\mathbf s)\,e^{\lambda\sum_i\hat g\cdot\mathbf s_i},\qquad G=F-\lambda m\\
K_A(t)=K_F(t+\lambda)-K_F(\lambda)\ \Rightarrow\ \kappa^A_n=K_F^{(n)}(\lambda)\\
\Delta\mu_i=\lambda\kappa_{2,i}+\tfrac{\lambda^2}{2}\kappa_{3,i},\qquad \kappa^A_{2,i}=\kappa_{2,i}+\lambda\kappa_{3,i}\\
\chi_{\rm swarm}=\chi_0/(1-\beta J_0\chi_0),\qquad g=\beta J_0\chi_0=1-V_{\rm indep}/V\\
\Delta\boldsymbol\mu=\lambda\,C_F\,\hat g
\end{gather*}\vspace{-8pt}}

\hobs{../figures/summary_obs.pdf}{(a) Each agent's push toward the assigned goal vs.\ its free-week fluctuation along $\hat g$, both divided by the pair mean. The tilt predicts the solid line, uniform translation the dashed one; the regime-I pairs slope the wrong way. (b) Change in variance along $\hat g$ (colored), along transverse directions (bars), and for an exact tilt of a bounded mean-field swarm at the same push ($\times$, synthetic). The goal disperses the swarm along $\hat g$ instead of narrowing it.}

\htelos{Mostly a negative result, useful only in a narrow sense. The appealing promise, reading how steerable a swarm is from how it fluctuates when nobody steers it, does not hold here. Goal changes are huge kicks (median about 2.5 SDs), not small fields, and the agents that wander most are not the ones that follow most. What an operator learns: an assigned goal acts like a target everyone converges on and a new mode the swarm toggles collectively, not a gentle tilt of existing preferences. So pre-goal behavior is a poor forecast of compliance. The instrument (goal-direction projection, noise-corrected variance) is cheap on any agent log, but the fluctuation--response physics it was built to test is, in this village, a curiosity.}

\hbottom{In the AI Village, goals are quenches toward a common target, not Legendre pushes: free-week fluctuations do not predict who follows a goal, and the variance moves the wrong way.}
